把上亿条记录散进 64 个分片，最省事的做法是拿记录编号除以 64 取余数。这行代码平常到没人追问一句：为什么散完之后，仍然可以把两个编号之和直接算在分片号上——先相加再取余，与先取余再相加，凭什么落在同一个分片？

取余确实丢信息，但丢得很讲究。整数按除以 \(m\) 的余数归成 \(m\) 类，加法和乘法可以原封不动地搬到类上：两类相加的结果，不依赖你从每一类里挑了哪个数当代表。无限多的整数被压进有限个抽屉，两种运算毫发无损地跟了过来。本章的主线就在这里，它也是\hyperref[sec:01-Mathematical-Language/02-Sets-Relations-Functions/02]{``关系与等价类''}中那套商结构第一次在具体运算上兑现。

代价落在除法头上。有理数里除零之外处处能除，抽屉里不行：模 12 时 \(3\times4\) 落回 0，两个非零的类乘出零，谁也约不掉谁。能约还是不能约，由一个此前只出现在算法里的量说了算——\(a\) 在模 \(m\) 下可逆，当且仅当 \(\gcd(a,m)=1\)。Euclid 算法与代数结构就在这一条上焊死：Bézout 恒等式既判定逆元是否存在，又把它算出来。

可逆元一旦立住，就能问指数。与 \(n\) 互素的余数类在乘法下自成一个只有 \(\varphi(n)\) 个元素的有限集合，反复自乘必然绕回起点，Fermat 与 Euler 定理把绕回的步数钉在 \(\varphi(n)\) 以内。指数于是可以先化简再计算，天文数字的幂变成几步平方，公钥密码的地基正是这块。最后一步是拆：模数两两互素时，``分别在几个小模下求解再拼装''与``在一个大模下求解''给出同一个答案，且拼法唯一。中国剩余定理讲的就是这件事，它是分片、并行与分治在数论里的原型。

\begin{center}
\begin{tikzpicture}[x=1cm,y=1cm,font=\small,>=stealth]
  \foreach \x/\a/\b in {0/{整除与 gcd}/{Euclid 变成算法}, 3.05/{余数成类}/{加乘搬得过去}, 6.1/{可逆元有限}/{指数出现周期}, 9.15/{模数互素}/{拆开算再拼回}}{
    \draw[black!45,fill=black!4] (\x,-0.8) rectangle (\x+2.5,0.8);
    \node[align=center] at (\x+1.25,0) {\a\\\b};}
  \foreach \x in {2.5,5.55,8.6}{\draw[->,black!55] (\x,0) -- (\x+0.5,0);}
  \foreach \x/\s in {0/{第一节}, 3.05/{第二节}, 6.1/{第三节}, 9.15/{第四节}}{
    \node[black!55,font=\footnotesize] at (\x+1.25,-1.2) {\s};}
\end{tikzpicture}

\emph{四节沿同一条线索推进：先把公因数变成可高效计算的东西，再把余数变成可运算的对象，然后追踪可逆元自乘的轨迹，最后把一个大模数上的问题拆成互不干扰的小问题。}
\end{center}

\hyperref[sec:04-Discrete-Mathematics/04-Number-Theory/01]{``整除、最大公约数与 Euclid 算法''}从一间要铺整砖的房间进入。反复从矩形里切掉最大的正方形，就是辗转相除；每切一轮，公因数集合原样不变。这一节的收获不止于算得快：Bézout 恒等式把``\(a\) 与 \(b\) 互素''翻译成一个能参与加减乘的等式 \(ax+by=1\)，后面三节几乎每一处关键推理都要调用这次翻译。

\hyperref[sec:04-Discrete-Mathematics/04-Number-Theory/02]{``模运算与模逆''}把余数升格为运算对象。同余是等价关系，加法和乘法在等价类上有定义，于是可以算到一半就取余，中间结果永远压在小范围内。消去律却失守了，\(2\cdot1\equiv2\cdot4\pmod 6\) 而 \(1\not\equiv4\)。这一节给出失守的判据，也给出可逆时的构造方法。

\hyperref[sec:04-Discrete-Mathematics/04-Number-Theory/03]{``Fermat 与 Euler 定理''}追踪 \(a,a^2,a^3,\dots\) 在模 \(n\) 下的轨迹。乘以一个可逆元只是把互素余数类重新排座位，排来排去总要回到出发点，回程不会超过 \(\varphi(n)\) 步。指数因此可以先模 \(\varphi(n)\) 约化。这一节也说清了定理的边界：\(\varphi(n)\) 是周期的上界而非最小周期，底数不互素时结论直接失效。

\hyperref[sec:04-Discrete-Mathematics/04-Number-Theory/04]{``中国剩余定理''}处理反方向的问题：几个局部余数能不能唯一确定一个全局整数。模数两两互素时能，而且对应关系同时保住加法和乘法，大模数上的计算可以拆到小模数上并行完成再拼回。互素性塌掉时，问题从``怎样构造''退化成``是否相容''。

四节按依赖顺序排列，顺读最省力。若手头有具体问题，也可以直接切入：要判断一个整数方程可不可解，或要算出模逆，进第一、二节；要处理巨大的指数或读懂 RSA 的骨架，进第三节；要把大模数上的运算拆开或把多路余数拼回，进第四节。若只是从纠错码那边回补基础，读完模逆一篇就足以进入\hyperref[ch:031]{``编码与多项式方法：用结构化冗余对抗丢失和错误''}。

读完之后可检验的变化是一种反射：写下 \(a^{-1}\) 之前先问 \(\gcd(a,m)\) 是多少，写下``指数模 \(\varphi(n)\)''之前先问底数与模数互不互素。这两个检查点在密码学、哈希设计和有限域计算里会反复出现，漏掉其中一个，得到的往往不是错误提示，而是一个安静地跑出错误结果的程序。
